% preamble.tex — packages and shared macros (coordinator edits only)
% Chapter agents: do NOT edit this file unless asked.
% wersja robocza 3: numerowane przypisy (biblatex numeric + hyperref),
% spokojniejszy skład (interlinia, odstęp akapitów, linia przypisów).

\documentclass[11pt,a4paper,twoside,openright]{book}

% Shared title/author macros (must load before hyperref metadata)
% metadata.tex — shared title/author metadata (coordinator edits only)
% Do not put chapter prose here.

\newcommand{\ThesisTitle}{%
IDA: projekt eksperymentalnej platformy badawczej do personalizacji
koncepcji wnętrz na podstawie preferencji estetycznych wspomaganej
sztuczną inteligencją}

\newcommand{\ThesisTitleShort}{IDA: platforma badawcza personalizacji wnętrz}

\newcommand{\ThesisAuthor}{Jakub Palka}
\newcommand{\ThesisSupervisor}{prof.\ dr hab.\ Andreas Guskos}
\newcommand{\ThesisAffiliation}{Akademia Sztuki w Szczecinie}
\newcommand{\ThesisDiscipline}{Sztuki plastyczne i konserwacja dzieł sztuki}
\newcommand{\ThesisYear}{2026}
\newcommand{\ThesisVersion}{robocza-3}
\newcommand{\ThesisVersionLabel}{wersja robocza 3}


% --- Encoding / engine (LuaLaTeX) ---
\usepackage{fontspec}
\usepackage{microtype}
\setmainfont{Latin Modern Roman}
\setsansfont{Latin Modern Sans}
\setmonofont{Latin Modern Mono}

% --- Language (babel + LuaLaTeX; avoids polyglossia/cleveref clash) ---
\usepackage[english,main=polish]{babel}

% --- Page geometry ---
\usepackage[
  a4paper,
  margin=2.5cm,
  bindingoffset=0.5cm
]{geometry}

% --- Typography: oddech tekstu ---
\usepackage{setspace}
\setstretch{1.16}
\usepackage{indentfirst}
\setlength{\parindent}{1.4em}
\setlength{\parskip}{0.32em plus 0.12em minus 0.05em}
\usepackage{csquotes}

% Spisy treści bez dodatkowego odstępu akapitów
\let\origtableofcontents\tableofcontents
\renewcommand{\tableofcontents}{{\setlength{\parskip}{0pt}\origtableofcontents}}
\let\origlistoffigures\listoffigures
\renewcommand{\listoffigures}{{\setlength{\parskip}{0pt}\origlistoffigures}}
\let\origlistoftables\listoftables
\renewcommand{\listoftables}{{\setlength{\parskip}{0pt}\origlistoftables}}

% --- Przypisy: czytelny rozmiar, cienka krótka linia ---
\renewcommand{\footnoterule}{%
  \kern 6pt
  \hrule width 0.22\columnwidth height 0.3pt
  \kern 4pt}
\setlength{\skip\footins}{14pt plus 4pt minus 2pt}
\setlength{\footnotesep}{7pt}
\usepackage{xcolor}
\definecolor{muted}{HTML}{3F5A73}
\definecolor{mutedgray}{HTML}{6B6B6B}

% --- Math / units (optional light use) ---
\usepackage{amsmath,amssymb}

% --- Tables / figures ---
\usepackage{booktabs}
\usepackage{tabularx}
\usepackage{longtable}
\usepackage{multirow}
\usepackage{graphicx}
\graphicspath{{figures/}{../figures/}}
\usepackage{pgfplots}
\pgfplotsset{compat=1.18}
\usepackage{float}
\usepackage{caption}
\usepackage{subcaption}
\captionsetup{
  font=small,
  labelfont=bf,
  format=hang
}

% --- Cross-refs (hyperref before cleveref): stonowane kolory, bez ramek ---
\usepackage{hyperref}
\hypersetup{
  colorlinks=true,
  linkcolor=muted,
  citecolor=muted,
  urlcolor=muted,
  filecolor=muted,
  bookmarksnumbered=true,
  pdftitle={\ThesisTitleShort},
  pdfauthor={\ThesisAuthor}
}
\usepackage[nameinlink]{cleveref}
% Polish cleveref names (Debian TeX Live may lack `polish` option)
\crefname{chapter}{rozdz.}{rozdz.}
\Crefname{chapter}{Rozdział}{Rozdziały}
\crefname{section}{podrozdz.}{podrozdz.}
\Crefname{section}{Podrozdział}{Podrozdziały}
\crefname{figure}{ryc.}{ryc.}
\Crefname{figure}{Rycina}{Ryciny}
\crefname{table}{tab.}{tab.}
\Crefname{table}{Tabela}{Tabele}
\crefname{appendix}{zał.}{zał.}
\Crefname{appendix}{Załącznik}{Załączniki}
\crefname{equation}{równ.}{równ.}
\Crefname{equation}{Równanie}{Równania}

% --- Lists ---
\usepackage{enumitem}
\setlist{nosep,leftmargin=*}

% --- Code / prompts (short listings only) ---
\usepackage{listings}
\lstset{
  basicstyle=\ttfamily\small,
  breaklines=true,
  frame=single,
  columns=fullflexible
}

% --- Bibliography: numerki w przypisach ---
% style=numeric: w przypisie jest numer pozycji, nie opis autor–rok w zdaniu.
% hyperref: numer w przypisie -> wpis w bibliografii; backref: wpis -> strony cytowania.
\usepackage[
  backend=biber,
  style=numeric,
  autocite=footnote,
  hyperref=true,
  backref=true,
  backrefstyle=three,
  maxbibnames=99,
  doi=true,
  url=true,
  isbn=false,
  giveninits=true,
  uniquename=init
]{biblatex}
\addbibresource{bibliography/references.bib}

% --- Draft helpers ---
\newcommand{\Todo}[1]{\textcolor{red}{\textbf{[TODO: #1]}}}
\newcommand{\AgentNote}[1]{\textcolor{teal}{\textit{[agent: #1]}}}
\newcommand{\NeedsCite}[1]{\textcolor{red}{[ŹRÓDŁO DO UZUPEŁNIENIA]}}
% Brak źródła dla twierdzenia spoza obserwacji autora:
\newcommand{\BrakZrodla}{\textcolor{red}{[ŹRÓDŁO DO UZUPEŁNIENIA]}}

% --- Label namespaces (multi-agent safety) ---
% Prefer: \label{ch02:sec:biophilia}, \label{fig:ch07:flow}, \label{tab:ch06:instruments}
\newcommand{\ChLabel}[2]{\label{ch#1:#2}}

% --- Chapter status banner: w wersji do czytania wyłączony (nagłówki STATUS zostają w plikach) ---
\newif\ifThesisDraft
\ThesisDraftfalse
\newcommand{\ChapterStatus}[2]{%
  \ifThesisDraft
    \begin{center}
      \footnotesize\color{gray}
      Status: \textbf{#1} \quad|\quad Owner: \texttt{#2} \quad|\quad \ThesisVersion
    \end{center}
    \vspace{0.5em}
  \fi
}

% --- Consistent hypothesis / RQ macros ---
\newcommand{\RQ}[1]{\textbf{RQ#1}}
\newcommand{\Hyp}[1]{\textbf{H#1}}
